\chapter{Conclusão}\label{chp:CONCLUSAO}

Este capítulo apresenta as considerações finais sobre o desenvolvimento do DataRural. Inicialmente, são retomados o objetivo do trabalho e os principais resultados alcançados com a construção da plataforma. Em seguida, são discutidas as limitações do estágio atual e as possibilidades de continuidade do projeto.

\section{Considerações finais}

Este trabalho teve como objetivo projetar e desenvolver uma plataforma institucional para apoiar a publicação, a organização, a consulta e o compartilhamento de datasets da UFRRJ. A proposta foi motivada pela necessidade de oferecer aos dados um tratamento que não se limitasse ao armazenamento de arquivos, incorporando informações descritivas, condições de acesso, autoria e histórico de atualizações.

Como resultado, foi desenvolvido o DataRural, uma aplicação web capaz de receber um ou mais arquivos CSV por versão e associá-los aos metadados definidos para o dataset. A plataforma disponibiliza consulta e download de conteúdos públicos sem exigir autenticação, ao mesmo tempo que reserva as operações de publicação e gerenciamento aos usuários identificados. A busca, os filtros por área temática e a prévia tabular complementam esse fluxo, permitindo observar o contexto e parte da estrutura do material antes de obtê-lo.

O versionamento constitui um dos resultados centrais da implementação. Em vez de substituir os arquivos de uma publicação, cada atualização pode ser registrada separadamente, mantendo a relação com o dataset original. A exclusão lógica e a possibilidade de restauração complementam esse mecanismo, enquanto os registros de auditoria documentam alterações realizadas nas entidades da aplicação. Esses recursos concretizam os requisitos de preservação do histórico apresentados no Capítulo~\ref{chp:PROPOSTA}.

A colaboração foi incorporada por meio de grupos e papéis contextuais. Essa organização permite que um dataset seja mantido por uma equipe sem conceder as mesmas permissões a todos os participantes. A separação entre datasets públicos e privados também possibilita utilizar a plataforma tanto para divulgação quanto para atividades internas. As verificações de acesso são executadas no servidor e utilizam a propriedade do conjunto, sua associação a um grupo e o papel exercido pelo usuário.

Do ponto de vista técnico, a aplicação foi organizada em módulos correspondentes aos principais domínios do sistema. A integração entre AdonisJS, React e Inertia.js permitiu manter o processamento das requisições e das permissões no servidor, oferecendo uma interface composta por elementos reutilizáveis. O PostgreSQL concentra os dados relacionais, enquanto os anexos permanecem separados e são associados às respectivas versões. A implementação descrita no Capítulo~\ref{chp:EXPERIMENTOS} fornece, portanto, uma base funcional para a continuidade do DataRural.

O trabalho não procura substituir repositórios científicos de alcance amplo ou plataformas comunitárias já consolidadas. Sua contribuição está na reunião de recursos de publicação, controle de acesso, colaboração e acompanhamento de versões em uma solução direcionada ao contexto institucional apresentado. Com isso, o projeto demonstra como os requisitos levantados podem ser transformados em um sistema único, mantendo pontos de extensão para necessidades que ultrapassam o escopo desta etapa.

\section{Limitações e trabalhos futuros}

O estágio atual do DataRural apresenta limitações que orientam possíveis trabalhos futuros. A primeira está relacionada aos formatos aceitos. A publicação foi implementada inicialmente para arquivos CSV, com limite de 25~MB por arquivo. Embora esse recorte atenda ao objetivo definido para o projeto, ele não contempla outros materiais que podem acompanhar uma pesquisa, como planilhas, documentos geoespaciais, imagens ou formatos colunares. Uma extensão da plataforma pode criar validadores e visualizadores específicos para novos formatos, preservando o modelo de metadados e versões já adotado.

O tratamento do próprio CSV também pode ser ampliado. A visualização atual apresenta uma amostra tabular e o indicador de usabilidade considera aspectos estruturais observados durante o envio. Esses recursos não verificam a validade científica dos valores nem substituem a documentação produzida pelo responsável. Como continuidade, podem ser acrescentados dicionários de dados, identificação configurável de delimitadores e codificações, regras de validação por coluna e comparações entre versões. Essas funcionalidades permitiriam informar com maior detalhe como a estrutura do conjunto se altera ao longo do tempo.

Outra limitação é o armazenamento dos anexos no sistema de arquivos local. A camada utilizada pela aplicação permite que essa estratégia seja substituída posteriormente sem alterar o modelo principal dos datasets. Um trabalho futuro pode integrar um serviço de armazenamento de objetos e definir procedimentos de cópia de segurança, recuperação e retenção dos arquivos. Também pode ser estabelecida uma política para remoção definitiva de versões marcadas como excluídas, respeitando as necessidades administrativas e as restrições aplicáveis a cada conjunto.

A implementação priorizou os fluxos funcionais da plataforma e não incluiu, no escopo apresentado, uma avaliação sistemática com usuários da universidade. Dessa forma, não é possível concluir, a partir deste trabalho, como diferentes públicos percebem a facilidade de uso ou quais metadados consideram mais importantes durante a publicação e a busca. Estudos futuros podem conduzir testes de usabilidade com pesquisadores, estudantes e servidores, utilizando os resultados para revisar os formulários, a navegação, os papéis dos grupos e a apresentação das informações.

Também são necessárias avaliações técnicas antes de uma disponibilização institucional em maior escala. Testes de carga podem observar o comportamento da aplicação com volumes superiores de arquivos e acessos simultâneos, enquanto testes de segurança podem examinar os fluxos de autenticação, autorização e download. Testes automatizados também podem ser acrescentados para verificar as regras de negócio, especialmente nos cenários que combinam propriedade, visibilidade e papéis de grupo.

Por fim, a evolução do DataRural pode incluir integração com serviços institucionais de identidade, suporte completo a internacionalização e mecanismos de identificação persistente para publicações que necessitem de citação formal. Integrações com outros repositórios e catálogos também podem evitar que o acervo permaneça isolado. Essas extensões devem manter a separação de responsabilidades da arquitetura atual e ser acompanhadas por definições institucionais sobre curadoria, privacidade, licenciamento e conservação dos dados.
